% Báo cáo Lab 01 - Linear Regression (CSC14005). Biên dịch bằng XeLaTeX.
% Khuyến nghị 10-20 trang. Mọi số liệu lấy từ models/*.csv; mọi trích dẫn có trong references.bib.
\documentclass[12pt,a4paper]{article}
\usepackage[a4paper,margin=2.5cm]{geometry}
\usepackage{fontspec}
\setmainfont{TeX Gyre Termes}
\usepackage{amsmath,amssymb}
\usepackage{graphicx}
\usepackage{booktabs}
\usepackage{float}
\usepackage{caption}
\usepackage[numbers]{natbib}
\usepackage{hyperref}
\graphicspath{{figures/}}

\renewcommand{\contentsname}{Mục lục}
\renewcommand{\listfigurename}{Danh mục hình}
\renewcommand{\listtablename}{Danh mục bảng}
\renewcommand{\refname}{Tài liệu tham khảo}
\renewcommand{\figurename}{Hình}
\renewcommand{\tablename}{Bảng}

\begin{document}

\begin{titlepage}
  \centering
  {\large ĐẠI HỌC QUỐC GIA TP.HCM\\TRƯỜNG ĐẠI HỌC KHOA HỌC TỰ NHIÊN\\KHOA CÔNG NGHỆ THÔNG TIN\par}
  \vspace{3cm}
  {\LARGE\bfseries BÁO CÁO ĐỒ ÁN LAB 01\par}
  \vspace{0.5cm}
  {\Large Đề tài: % TODO(team): tên bài toán
  \par}
  \vspace{0.5cm}
  {\large Môn học: CSC14005 -- Nhập môn học máy\par}
  \vfill
  \begin{flushleft}
    \textbf{Nhóm:} % TODO(team): mã nhóm
    \\[0.3cm]
    \textbf{Sinh viên thực hiện:}\\
    % TODO(team): Họ tên -- MSSV
  \end{flushleft}
  \vfill
  {\today\par}
\end{titlepage}

\tableofcontents
\listoftables
\listoffigures
\clearpage

\section{Giới thiệu}
\subsection{Phân công công việc}
% TODO(team): bảng thành viên - nhiệm vụ - mô hình phụ trách - mức độ hoàn thành.

\subsection{Sử dụng công cụ AI}
% TODO(team): nêu rõ công cụ AI đã dùng, cho việc gì, và cách nhóm kiểm tra lại nội dung.

\section{Bài toán và động lực}
% 3.1: động lực chọn bài toán; mục tiêu dự đoán; biến mục tiêu; các biến đầu vào.

\section{Dữ liệu}
\subsection{Nguồn và mô tả dữ liệu}
% Nguồn (link), giấy phép, số mẫu, bảng mô tả các biến.
\subsection{Phân tích khám phá dữ liệu}
% Phân phối biến mục tiêu, tương quan, ngoại lệ, đa cộng tuyến. Hình từ notebooks/01_eda.ipynb.
\subsection{Tiền xử lý}
% Xử lý thiếu, mã hoá, chuẩn hoá, biến đổi; chia train/test (80/20, seed 42).

\section{Thiết kế mô hình}
\subsection{Cơ sở lý thuyết}
% OLS, Ridge, Lasso, Elastic Net, hàm cơ sở đa thức, biến đổi log mục tiêu, Huber.
\subsection{Các mô hình xây dựng}
% Mỗi mô hình: biến sử dụng, biến đổi, siêu tham số, lý do lựa chọn (yêu cầu 3.2).

\section{Đánh giá và so sánh}
\subsection{Chỉ số đánh giá}
% MAE, MSE, RMSE, R^2, cross-validation.
\subsection{Kết quả}
% Bảng từ models/metrics.csv; phân tích phần dư.
\subsection{Diễn giải hệ số}
% Chỉ dựa trên kết quả mô hình (yêu cầu 3.3), không dùng kiến thức miền bên ngoài.

\section{Ứng dụng}
% Tính năng, giao diện (ảnh chụp màn hình), đầu vào/đầu ra, cách tích hợp mô hình đã huấn luyện, link video demo.

\section{Kết luận}
% Kết quả chính, hạn chế, hướng cải tiến.

\nocite{*}
\bibliographystyle{plainnat}
\bibliography{references}

\end{document}
